\documentclass[12pt,a4paper]{article}

\usepackage[english]{babel}
\usepackage{geometry}
\usepackage{hyperref}
\usepackage{xcolor}
\usepackage{graphicx}
\usepackage{array}
\usepackage{booktabs}
\usepackage{fancyhdr}
\usepackage{listings}
\usepackage{amssymb}
\usepackage{stix2}
\usepackage{enumitem}

\geometry{margin=1in}

\pagestyle{fancy}
\fancyhf{}
\rhead{CS486 -- Group 06}
\lhead{Group Report}
\cfoot{\thepage}

\definecolor{codebg}{gray}{0.95}
\definecolor{codeframe}{gray}{0.7}

\lstset{
  basicstyle=\ttfamily\small,
  backgroundcolor=\color{codebg},
  frame=single,
  rulecolor=\color{codeframe},
  breaklines=true,
  breakatwhitespace=true,
  columns=fullflexible,
  keepspaces=true,
  literate={→}{$\rightarrow$}1 {←}{$\leftarrow$}1,
}

\title{\textbf{CS486 -- Campus Space Management System} \\ \vspace{0.5cm} \Large \textbf{Group Agents Report}}
\author{Group Number: 06 \\Nguyen Hong Tan Tai -- 24125078 \\Quach Thien Lac -- 24125092 \\ Nguyen Dinh Thien Loc -- 24125093 \\ Tran Ton Minh Ky -- 24125102}

\begin{document}

\maketitle

\begin{abstract}
This report documents the work of Group G06 on the CS486 database design agent project. Part~I covers the improvements made to the baseline 2-step teaching demo, extending it into a comprehensive 7-step pipeline with state tracking, validation, and explicit traceability. Part~II details how the markdown configuration files (\texttt{AGENTS.md}, \texttt{MEMORY.md}, \texttt{PROMPTS.md}, \texttt{SKILL.md}) were designed and structured to encode the teacher's guidelines as machine-readable agent rules. All work was performed through Antigravity CLI using multiple LLM models.
\end{abstract}


\newpage

\tableofcontents

\newpage

%% ============================================================
%%  PART I — AGENT IMPROVEMENTS
%% ============================================================

\part{Agent Improvements}

\section{LLM Models Used}

All agent work was performed through Antigravity CLI, using the following models:

\begin{itemize}
    \item \textbf{Claude Sonnet (Anthropic)} -- primary model for iterative development, prompt refinement, and file editing
    \item \textbf{Claude Opus (Anthropic)} -- used for complex reasoning tasks such as design validation and report generation
    \item \textbf{Gemini Flash (Google)} -- used for fast, low-cost drafts and quick syntax checks
    \item \textbf{Gemini Pro (Google)} -- used for cross-validation and alternative design perspectives
\end{itemize}

This multi-model approach enabled cross-validation of design decisions and provided diverse technical perspectives on implementation correctness.

\section{Agent Improvement Overview}

\subsection{Baseline Assessment}

The original teaching demo agent (hereafter called the \textit{baseline agent}) implemented a minimal 2-step workflow:

\begin{enumerate}
    \item Analyze business requirements
    \item Produce conceptual ERD using Chen's notation
\end{enumerate}

\noindent Output locations:
\begin{itemize}
    \item \texttt{docs/01-business-requirement-analysis.md}
    \item \texttt{docs/02-conceptual-design-erd.md}
\end{itemize}

\noindent Limitations identified:
\begin{itemize}
    \item No logical schema design (relational normalization)
    \item No design validation step (Normal Forms, constraint verification)
    \item No SQL database implementation (DDL)
    \item No sample data preparation (DML)
    \item No query design or business query support
    \item No state tracking across sessions
    \item No explicit prerequisite enforcement
    \item Implicit assumptions rather than explicit escalation
    \item No formalized ID naming conventions
\end{itemize}

\subsection{Performance Evaluation Methodology}

The group evaluated agent performance against these grading criteria at each refinement iteration:

\begin{table}[h]
\centering
\begin{tabular}{|p{3cm}|p{4cm}|p{4cm}|}
\hline
\textbf{Criterion} & \textbf{Baseline Status} & \textbf{Evaluation Method} \\
\hline
Completeness & Incomplete (5 of 7 steps missing) & Check all required deliverables exist \\
\hline
Traceability & Implicit, undocumented & Trace every table/constraint to requirement \\
\hline
Correctness & No validation mechanism & Verify Normal Forms and constraint consistency \\
\hline
Automation & No state tracking & Can agent resume without re-work? \\
\hline
Usability & Rules scattered & Are prerequisites and rules explicit? \\
\hline
\end{tabular}
\end{table}

\newpage

\section{Major Improvements}

\subsection{Improvement 1: Extended Pipeline (2 Steps → 7 Steps)}

\subsubsection{Problem}
The baseline agent produced only ERD; no database implementation artifacts were generated.

\subsubsection{Solution}
Extended the pipeline from 2 steps to a complete 7-step workflow:

\begin{enumerate}
    \item Business requirement analysis
    \item Conceptual ERD design (Crow's Foot notation)
    \item \textbf{Logical design} (relational schema, keys, normalization)
    \item \textbf{Design validation} (Normal Forms, constraint verification, traceability matrix)
    \item \textbf{Database implementation (DDL)} -- SQL \texttt{CREATE TABLE} with constraints, checks, defaults
    \item \textbf{Sample data preparation (DML)} -- \texttt{INSERT} statements for realistic testing
    \item \textbf{Query design (SQL)} -- 5+ business-critical queries with explanations
\end{enumerate}

\subsubsection{Output Files}
All artifacts in \texttt{outputs/}:
\begin{itemize}
    \item \texttt{01-business-req-analysis-G06.md}
    \item \texttt{02-erd-design-G06.md}
    \item \texttt{03-logical-design-G06.md}
    \item \texttt{04-design-validation-G06.md}
    \item \texttt{05-db-definition-G06.sql}
    \item \texttt{06-sample-data-G06.sql}
    \item \texttt{07-query-design-G06.sql}
\end{itemize}

\subsubsection{Impact}
\begin{itemize}
    \item All 7 required pipeline stages now produce formal, graded deliverables
    \item Complete path from requirement to executable SQL
    \item Business queries demonstrate system correctness
\end{itemize}

\subsection{Improvement 2: State Tracking with MEMORY.md}

\subsubsection{Problem}
No mechanism existed to track which step was done, in-progress, or blocked. Sessions risked re-doing prior work or forgetting locked decisions.

\subsubsection{Solution}
Introduced \texttt{MEMORY.md} as a centralized project state file (see Part~II, Section~\ref{sec:memory} for detailed structure).

\subsubsection{Impact}
\begin{itemize}
    \item Sessions can resume without re-reading entire prior outputs
    \item No repeated work; every decision is auditable
    \item Open questions are explicit, not silently resolved
    \item Reduced token usage per session
\end{itemize}

\subsection{Improvement 3: Detailed Pipeline Skill (SKILL.md)}

\subsubsection{Problem}
Agent workflow rules were scattered across comments in \texttt{AGENTS.md}. Prerequisites, validation criteria, and traceability rules were implicit and hard to enforce.

\subsubsection{Solution}
Created \texttt{.opencode/skills/db-design-pipeline/SKILL.md} (synced to \texttt{.openclaw/} for tool portability):

\begin{itemize}
    \item \textbf{Step-by-step templates}: Each of 7 steps has explicit input, output template, "done" criteria
    \item \textbf{Prerequisite checks}: Step $N$ cannot run if step $N-1$ is incomplete
    \item \textbf{Validation criteria}: What makes output "done" vs "needs revision"
    \item \textbf{Traceability rules}: Every table/constraint must cite requirement line number
    \item \textbf{Token discipline}: Only regenerate files user explicitly names
\end{itemize}

\subsubsection{Impact}
\begin{itemize}
    \item Agent now enforces strict workflow order
    \item Prevents out-of-sequence work (e.g., cannot skip to queries before validation)
    \item Clear, auditable rules for what constitutes acceptable output
\end{itemize}

\subsection{Improvement 4: Explicit Traceability Matrix}

\subsubsection{Problem}
Baseline agent recorded entities and relationships but did not explicitly map them back to requirement sentences.

\subsubsection{Solution}
Added comprehensive traceability matrix in step 1 analysis and MEMORY.md:

\begin{small}
\begin{verbatim}
Entity/Requirement Traceability:
- User entity ← req: user ID, full name, email, phone, 
                    role, department, account status
- Space entity ← req: space code, name, type, 
                      building/floor/room, capacity, status, policy
- Facility entity ← req: facility ID, facility name
- Space_Facility junction ← req: list of facilities per space
- BookingRequest ← req: booking with space, times, 
                        purpose, participants, status
- BookingApproval ← req: approval/rejection by staff
- Maintains (maintenance) ← req: maintenance records per space
- Lookup tables: UserRole, RoomStatus, RequestStatus, 
                 MaintenanceStatus
\end{verbatim}
\end{small}

\subsubsection{Impact}
\begin{itemize}
    \item Every design element has explicit justification from requirement
    \item Grader can verify no invented features
    \item Supports systematic completeness validation
\end{itemize}

\subsection{Improvement 5: Naming Conventions \& ID Standardization}

\subsubsection{Problem}
Business requirement was vague about ID formats and naming schemes. Agents were inventing ad-hoc conventions.

\subsubsection{Solution}
Locked explicit naming standards in MEMORY.md:

\begin{table}[h]
\centering
\begin{tabular}{|l|p{5cm}|}
\hline
\textbf{Entity} & \textbf{ID Format} \\
\hline
Space & 3--5 letter codes (e.g., \texttt{HT1}, \texttt{I12b}) \\
\hline
User & 8-digit numeric, generated on creation \\
\hline
Facility & 6-char: 3 alpha + 3 numeric (e.g., \texttt{chr055}) \\
\hline
Booking Request & 8-char lowercase alphanumeric (non-enumerable) \\
\hline
Maintenance & 6-char lowercase alphanumeric \\
\hline
Full Name & Computed: \texttt{first\_name + ' ' + last\_name} \\
\hline
Space Location & Computed: \texttt{building + floor + room\_number} \\
\hline
\end{tabular}
\end{table}

\subsubsection{Impact}
\begin{itemize}
    \item Consistent, auditable naming across all design phases
    \item Avoids design thrashing from convention changes
    \item Supports requirement traceability requirements
\end{itemize}

\subsection{Improvement 6: Validation Step (Step 4)}

\subsubsection{Problem}
ERD could be syntactically correct but violate business rules or fail normalization checks. No formal validation before SQL generation.

\subsubsection{Solution}
Inserted explicit \textbf{Design Validation} step (step 4):

\begin{itemize}
    \item Verify logical schema against ERD (all relationships properly represented)
    \item Check Normal Forms (3NF or BCNF with justification)
    \item Validate constraint coverage (no business rules omitted)
    \item Trace every table and column back to requirement
    \item Flag forward-references, circular dependencies, missing links
\end{itemize}

\subsubsection{Impact}
\begin{itemize}
    \item Early detection of design flaws before SQL generation
    \item Prevents invalid databases from being deployed
    \item Formal quality assurance gate before implementation step
\end{itemize}

\subsection{Improvement 7: Prerequisite Enforcement}

\subsubsection{Problem}
Agent could jump to any step (e.g., write queries without completing validation). No prevention of out-of-order work.

\subsubsection{Solution}
Formalized prerequisite rules in SKILL.md:

\begin{itemize}
    \item Step 2 (ERD) requires step 1 (analysis) $\checkmark$ done
    \item Step 3 (logical) requires steps 1--2 $\checkmark$ done
    \item Step 4 (validation) requires steps 1--3 $\checkmark$ done
    \item Step 5 (DDL) requires steps 1--4 $\checkmark$ done
    \item Step 6 (sample data) requires steps 1--5 $\checkmark$ done
    \item Step 7 (queries) requires steps 1--6 $\checkmark$ done
    \item \textbf{Override}: User can explicitly say \texttt{skip ahead anyway} only with documented justification
\end{itemize}

\subsubsection{Impact}
\begin{itemize}
    \item Prevents invalid or incomplete submissions
    \item Enforces proper design discipline and methodology
    \item Provides clear feedback when prerequisites are unmet
\end{itemize}

\subsection{Improvement 8: Multi-Tool Support (PROMPTS.md, setup.md)}

\subsubsection{Problem}
Baseline agent assumed single tool (OpenCode). No guidance for other tools (Antigravity, Codex CLI, OpenClaw).

\subsubsection{Solution}
Created infrastructure files:

\begin{itemize}
    \item \textbf{PROMPTS.md}: Five templated prompts for different scenarios (setup, session start, continue, fix one file, answer open question)
    \item \textbf{md/setup.md}: Per-tool installation and launch instructions
\end{itemize}

\subsubsection{Impact}
\begin{itemize}
    \item Same workflow functions consistently across multiple platforms
    \item Reduced setup friction and learning curve for team members
    \item Design rules remain portable and tool-independent
\end{itemize}

\newpage

\section{Quantitative Comparison}

\begin{table}[h]
\centering
\begin{tabular}{|l|c|c|}
\hline
\textbf{Metric} & \textbf{Baseline Agent} & \textbf{Improved Agent} \\
\hline
Pipeline steps & 2 & 7 \\
\hline
Output deliverables & 2 & 7 \\
\hline
State tracking & None & MEMORY.md \\
\hline
Prerequisite enforcement & No & Yes \\
\hline
Validation step & No & Yes \\
\hline
Traceability documentation & Implicit & Explicit matrix \\
\hline
Naming standards & Vague & Locked (8 standards) \\
\hline
Open question handling & Silent assumptions & Explicit escalation \\
\hline
Tool support & 1 (OpenCode) & 4 (+ PROMPTS templates) \\
\hline
Infrastructure files & 1 (AGENTS.md) & 3 (AGENTS.md, MEMORY.md, PROMPTS.md) \\
\hline
Skill templates & Inline comments & 7 dedicated step templates \\
\hline
\end{tabular}
\end{table}

\section{Evaluation Methodology Per Refinement Cycle}

Each improvement was evaluated against the following process:

\subsection{Cycle 1: Extended Pipeline}
\begin{enumerate}
    \item \textbf{Baseline}: Run step 1 (analysis) on business requirement and review traceability
    \item \textbf{Assessment}: Missing steps 3--7; cannot generate complete database
    \item \textbf{Refinement}: Add steps 3--7 to AGENTS.md workflow order
    \item \textbf{Re-evaluation}: All 7 steps now produce outputs; compare against rubric checklist
    \item \textbf{Result}: Pass — all required deliverables present
\end{enumerate}

\subsection{Cycle 2: State Tracking}
\begin{enumerate}
    \item \textbf{Baseline}: Run step 2 (ERD) and end session; resume next session and agent re-reads all prior outputs
    \item \textbf{Assessment}: Inefficient workflow; excessive token expenditure; risk of lost context
    \item \textbf{Refinement}: Create MEMORY.md; add rule to AGENTS.md: ``always read MEMORY.md first''
    \item \textbf{Re-evaluation}: Next session: agent reads MEMORY.md, knows exact state, continues without re-work
    \item \textbf{Result}: Pass — significant token savings per session
\end{enumerate}

\subsection{Cycle 3: Validation Step}
\begin{enumerate}
    \item \textbf{Baseline}: Generate step 5 (DDL) directly from step 2 (ERD)
    \item \textbf{Assessment}: No mechanism to check Normal Forms or missing constraints; invalid SQL possible
    \item \textbf{Refinement}: Insert step 4 (validation); require it before DDL generation
    \item \textbf{Re-evaluation}: Step 4 output now traces every table back to requirement; catches design flaws early
    \item \textbf{Result}: Pass — design quality gate successfully implemented
\end{enumerate}

\subsection{Cycles 4--8: Remaining Improvements}
Similar iterative cycles applied for naming standards, traceability matrix, prerequisite enforcement, and multi-tool support. Each cycle follows this methodology:
\begin{enumerate}
    \item Identified a specific weakness in current rules
    \item Proposed targeted infrastructure or rule addition
    \item Updated AGENTS.md or created new supporting file
    \item Re-ran agent on example input to verify improvement
    \item Recorded result in MEMORY.md with rationale
\end{enumerate}

\newpage

%% ============================================================
%%  PART II — MARKDOWN FILE SETUP
%% ============================================================

\part{Markdown Configuration File Setup}

This part documents how the markdown files were designed and structured to encode the teacher's guidelines as machine-readable rules for AI agents.

\section{Teacher's Guidelines (Reference Project)}

The teacher's reference project (\texttt{reference/}) establishes the following requirements that our configuration files must satisfy:

\subsection{Project Goal (Section 2 of Reference)}

The agent must read the business requirement and generate seven ordered artifacts:

\begin{enumerate}
    \item Business Requirement Analysis
    \item Conceptual Database Design
    \item Logical Database Design
    \item Database Design Validation
    \item Database Implementation
    \item Sample Data Preparation
    \item Query Design
\end{enumerate}

The teacher's baseline agent only implemented steps 1--2. Our configuration files encode all seven.

\subsection{Required Output Artifacts (Section 6 of Reference)}

The teacher specifies exact output file naming with a \texttt{G<Group number>} placeholder. For Group G06, the actual filenames are:
\begin{lstlisting}
01-business-req-analysis-G06.md
02-erd-design-G06.md
03-logical-design-G06.md
04-design-validation-G06.md
05-db-definition-G06.sql
06-sample-data-G06.sql
07-query-design-G06.sql
\end{lstlisting}

These are encoded verbatim in \texttt{AGENTS.md} Section ``Required Outputs,'' with \texttt{G06} hardcoded so the agent never needs to infer the group number.

\subsection{Cost Control Guidelines (Section 7 of Reference)}

The teacher emphasizes token discipline:
\begin{itemize}
    \item Use a cheaper model for early drafts; a stronger model for validation
    \item Do not repeatedly regenerate all files from scratch
    \item Ask the agent to update only the specific file or section that needs improvement
    \item Keep prompts short, clear, and specific
\end{itemize}

These rules are encoded in \texttt{AGENTS.md} Section ``Token / cost discipline'' and operationalized through \texttt{MEMORY.md} (which avoids re-reading completed outputs) and \texttt{PROMPTS.md} (which provides minimal, targeted prompt templates).

\subsection{Design Rules (Section 8 of Reference)}

The teacher requires:
\begin{itemize}
    \item Record assumptions explicitly
    \item Record open questions explicitly
    \item Preserve traceability: requirement $\rightarrow$ entity $\rightarrow$ relationship $\rightarrow$ table $\rightarrow$ constraint
    \item Use Mermaid \texttt{erDiagram} syntax for ERD
    \item Do not silently invent business rules
\end{itemize}

All five rules are replicated verbatim in \texttt{AGENTS.md} Section ``Design Rules.''

\subsection{Academic Integrity (Section 8 of Reference)}

Students must be able to explain how the agent was configured, improved, and why the final design is valid. This report serves that purpose.

\section{File Architecture Overview}

\begin{table}[h]
\centering
\begin{tabular}{|l|p{4.5cm}|p{4.5cm}|}
\hline
\textbf{File} & \textbf{Purpose} & \textbf{Read by} \\
\hline
\texttt{AGENTS.md} & Global rules, workflow order, design principles, DBMS choice & Antigravity, OpenCode, Codex CLI (auto-read) \\
\hline
\texttt{MEMORY.md} & Live project state: pipeline progress, locked decisions, open questions & All tools (via AGENTS.md instruction) \\
\hline
\texttt{PROMPTS.md} & Reusable prompt templates for common workflows & Human operators \\
\hline
\texttt{SKILL.md} & Detailed per-step templates, prerequisites, validation criteria & OpenCode (native), OpenClaw (synced copy), Antigravity/Codex (via AGENTS.md pointer) \\
\hline
\texttt{md/setup.md} & Per-tool installation and launch instructions & Human operators \\
\hline
\texttt{md/pipeline.md} & Project structure tree and pipeline overview & Human operators \\
\hline
\end{tabular}
\caption{Configuration file architecture}
\end{table}

\subsection{Design Principle: Separation of Rules and State}

A key design decision is the separation between \texttt{AGENTS.md} (static rules) and \texttt{MEMORY.md} (dynamic state):

\begin{itemize}
    \item \textbf{AGENTS.md} contains rules that rarely change: workflow order, output file naming, DBMS choice, design rules, traceability requirements. An agent reads this once per session.
    \item \textbf{MEMORY.md} contains state that changes after every task: which steps are done, locked decisions, open questions, assumptions. This is the ``save file'' that enables cheap session resumption.
\end{itemize}

This separation follows the teacher's cost-control guideline: the agent reads \texttt{MEMORY.md} instead of re-reading all 7 output files, significantly reducing token usage per session.

\section{AGENTS.md: Structure and Constraints}

\texttt{AGENTS.md} is the primary configuration file. It is auto-read by Antigravity, OpenCode, and Codex CLI at the start of every session.

\subsection{Header and Tool Compatibility}

\begin{lstlisting}
> This file is the shared instruction set for every
> agent tool used on this project (Antigravity,
> OpenCode, Codex CLI, OpenClaw, and others that read
> AGENTS.md natively).
\end{lstlisting}

This header establishes \texttt{AGENTS.md} as the single source of truth for all tools.

\subsection{Recurring Context}

Provides constants the agent needs every session:
\begin{itemize}
    \item Root directory: \texttt{./}
    \item Group number: G06
    \item Academic project flag (prioritize quality over speed)
    \item Filesystem verification: \texttt{ls -la outputs/} before assuming file existence
\end{itemize}

\subsection{Skill Location Pointers}

The detailed pipeline skill (\texttt{SKILL.md}) is stored at a canonical location:
\begin{lstlisting}
.opencode/skills/db-design-pipeline/SKILL.md
\end{lstlisting}

This canonical copy is synced to \texttt{.openclaw/skills/} via \texttt{scripts/sync-skills.sh}. Antigravity and Codex CLI reach the canonical path directly via the \texttt{AGENTS.md} pointer.

\subsection{MEMORY.md Directive}

\texttt{AGENTS.md} contains an explicit directive:

\begin{lstlisting}
## Always read MEMORY.md first

Before doing anything else in a session, read MEMORY.md
in the project root.
\end{lstlisting}

This ensures state tracking is not optional --- every agent tool is forced to read the state file before proceeding.

\subsection{Workflow Order Constraint}

Directly encoding the teacher's 7-step pipeline:

\begin{enumerate}
    \item Business requirement analysis
    \item Conceptual design (ERD, Crow's Foot notation)
    \item Logical design (tables, keys, normalization)
    \item Design validation (normal forms, constraint check, traceability check)
    \item Database implementation (DDL)
    \item Sample data preparation (DML)
    \item Query design (SQL queries answering business questions)
\end{enumerate}

The constraint ``Never produce a later-stage artifact if an earlier-stage artifact is missing, marked `in-progress,' or marked `needs revision' in MEMORY.md'' enforces strict prerequisite ordering.

\subsection{Design Rules (from Teacher's Guidelines)}

Five rules are encoded, matching the teacher's reference \texttt{AGENTS.md}:

\begin{enumerate}
    \item Record assumptions explicitly in the output file \textbf{and} in \texttt{MEMORY.md}
    \item Record open questions explicitly; do not silently resolve by guessing
    \item Preserve traceability: every table and constraint must trace back to \texttt{req/business-requirement.md}
    \item Use Mermaid \texttt{erDiagram} syntax for all ER diagrams
    \item Do not silently invent business rules not present in the requirement
\end{enumerate}

\subsection{Token Discipline Rules}

Three cost-control rules, directly implementing the teacher's Section 7 guidelines:

\begin{enumerate}
    \item Default to editing ONLY the specific file(s) the user names
    \item When asked to ``continue,'' consult \texttt{MEMORY.md} for the next pending step
    \item Summarize long file contents into \texttt{MEMORY.md} once finalized
\end{enumerate}

\section{MEMORY.md: State Tracking Design}
\label{sec:memory}

\texttt{MEMORY.md} is the project's ``save state'' file. It is updated by the agent after every task completion.

\subsection{Sections and Their Purpose}

\begin{table}[h]
\centering
\begin{tabular}{|l|p{7cm}|}
\hline
\textbf{Section} & \textbf{Purpose} \\
\hline
Group info & Group number (G06) and DBMS choice (SQL Server) \\
\hline
Pipeline status table & 7-row table with step name, output file, status, and notes \\
\hline
Locked decisions & Confirmed ID formats, naming conventions, normalization choices \\
\hline
Open questions & Ambiguous requirements escalated to user instead of guessed \\
\hline
Assumptions recorded & Implicit assumptions made explicit for auditability \\
\hline
Entity traceability snapshot & Short entity $\leftarrow$ requirement mappings \\
\hline
Known issues & Bugs and fixes needed in next session \\
\hline
\end{tabular}
\caption{MEMORY.md sections}
\end{table}

\subsection{Status Values}

Four status values are used consistently:
\begin{itemize}[nosep]
    \item \texttt{not started} --- file does not exist or is a 0-byte placeholder
    \item \texttt{in-progress} --- partially completed, not yet reviewable
    \item \texttt{done} --- complete and reviewed; later steps may depend on it
    \item \texttt{needs revision} --- was done but a problem was found; blocks dependent steps
\end{itemize}

\subsection{How MEMORY.md Implements Teacher's Cost Control}

The teacher's guideline ``Do not repeatedly regenerate all files from scratch'' is operationalized as follows:

\begin{enumerate}
    \item Agent reads \texttt{MEMORY.md} (small file) instead of re-reading all 7 output files
    \item Agent knows exact state: which steps are done, which are blocked, which need revision
    \item No re-derivation of prior decisions; they are cached in ``Locked decisions''
    \item Open questions are explicit, preventing the agent from silently guessing answers differently across sessions
\end{enumerate}

\section{PROMPTS.md: Prompt Template Design}

\texttt{PROMPTS.md} provides five reusable prompt templates, each designed to minimize token usage while maximizing agent effectiveness:

\begin{table}[h]
\centering
\begin{tabular}{|c|p{3cm}|p{6cm}|}
\hline
\textbf{\#} & \textbf{Template Name} & \textbf{Purpose} \\
\hline
0 & One-time setup & Set group to G06, DBMS to SQL Server, produce step 1 \\
\hline
1 & Session starter & Read MEMORY.md, report status in 3--5 bullets, wait \\
\hline
2 & Continue pipeline & Generate next pending file only \\
\hline
3 & Fix one file & Edit only the named file; do not touch others \\
\hline
4 & Answer open question & Resolve a question from MEMORY.md and apply answer \\
\hline
\end{tabular}
\caption{Prompt templates in PROMPTS.md}
\end{table}

Each template follows the teacher's principle: ``Ask the agent to update only the specific file or section that needs improvement.'' No template asks the agent to regenerate the entire pipeline.

\section{SKILL.md: Pipeline Enforcement}

The \texttt{SKILL.md} file at \texttt{.opencode/skills/db-design-pipeline/} provides detailed per-step enforcement:

\subsection{Step 0: Pre-flight Checks (Every Invocation)}

Before any work, the agent must:
\begin{enumerate}
    \item Read \texttt{MEMORY.md} to determine current state
    \item Run \texttt{ls -la outputs/} to verify filesystem matches \texttt{MEMORY.md}
    \item Check prerequisites: if the requested step's prior steps are not \texttt{done}, stop and tell the user
    \item If user says ``continue'' without naming a step, pick the first non-\texttt{done} step
\end{enumerate}

\subsection{Per-Step Specification}

Each of the 7 steps specifies:
\begin{itemize}
    \item \textbf{Input}: which prior output(s) to read
    \item \textbf{Template}: which template file to follow
    \item \textbf{Done criteria}: what makes the output acceptable
    \item \textbf{MEMORY.md update}: what to write after completion
\end{itemize}

This structure ensures the agent cannot produce incomplete or incorrect output, directly supporting the teacher's grading criteria for completeness and traceability.

\section{Tool Support Architecture}

\begin{table}[h]
\centering
\begin{tabular}{|l|p{3cm}|p{5cm}|}
\hline
\textbf{Tool} & \textbf{How it finds rules} & \textbf{How it finds SKILL.md} \\
\hline
Antigravity & Reads \texttt{AGENTS.md} automatically & Via AGENTS.md pointer to \texttt{.opencode/skills/} \\
\hline
OpenCode & Reads \texttt{AGENTS.md} automatically & Native skill loader at \texttt{.opencode/skills/} \\
\hline
Codex CLI & Reads \texttt{AGENTS.md} automatically & Via AGENTS.md pointer to \texttt{.opencode/skills/} \\
\hline
OpenClaw & Reads skill at \texttt{.openclaw/skills/} & Synced copy via \texttt{scripts/sync-skills.sh} \\
\hline
\end{tabular}
\caption{Tool support matrix}
\end{table}

All tools converge on the same rules through \texttt{AGENTS.md}. The canonical skill copy at \texttt{.opencode/skills/db-design-pipeline/SKILL.md} is synced to \texttt{.openclaw/} for OpenClaw compatibility.

\section{How Constraints Follow Teacher's Guidelines}

\begin{table}[h]
\centering
\begin{tabular}{|p{4cm}|p{4cm}|p{4cm}|}
\hline
\textbf{Teacher's Guideline} & \textbf{Our Implementation} & \textbf{Location} \\
\hline
7-step pipeline & Encoded as workflow order with prerequisite enforcement & AGENTS.md \S Workflow Order; SKILL.md per-step specs \\
\hline
Exact output filenames & Verbatim in Required Outputs section & AGENTS.md \S Required Outputs \\
\hline
Do not regenerate all files & Token discipline rules; MEMORY.md caching & AGENTS.md \S Token discipline; MEMORY.md \\
\hline
Update only specific file & Single-file prompt templates & PROMPTS.md templates 2--4 \\
\hline
Record assumptions & Explicit rule in Design Rules & AGENTS.md \S Design Rules; MEMORY.md \S Assumptions \\
\hline
Record open questions & Explicit escalation rule & AGENTS.md \S Design Rules; MEMORY.md \S Open questions \\
\hline
Preserve traceability & Entity-requirement mapping & AGENTS.md \S Design Rules; MEMORY.md \S Traceability \\
\hline
Use Mermaid ERD & Syntax requirement & AGENTS.md \S Design Rules \\
\hline
No invented rules & Guard rule & AGENTS.md \S Design Rules \\
\hline
Academic integrity & Agent is configured, not blindly trusted & This report; PROMPTS.md; MEMORY.md auditability \\
\hline
\end{tabular}
\caption{Mapping teacher's guidelines to our implementation}
\end{table}

\newpage

%% ============================================================
%%  SHARED CLOSING SECTIONS
%% ============================================================

\part*{} % invisible part break — just for visual separation
\addcontentsline{toc}{part}{Design Decisions, Limitations \& Conclusion}

\section{Design Decisions \& Lessons Learned}

\subsection{Why MEMORY.md Is Separate from AGENTS.md}

\begin{itemize}
    \item \textbf{AGENTS.md}: Global rules that apply to all projects (workflow order, DBMS choice, naming, design rules). Rarely changes.
    \item \textbf{MEMORY.md}: Project-specific state (current step, locked decisions, open questions, assumptions). Updates after every task completion.
\end{itemize}

This separation enables:
\begin{itemize}
    \item Reuse of AGENTS.md across future projects (just copy the template)
    \item Cheap session resumption (read MEMORY.md, not all 7 outputs)
    \item Explicit separation of rules vs. state
\end{itemize}

\subsection{Why Prerequisite Enforcement Matters}

Without prerequisites, agent could:
\begin{itemize}
    \item Generate SQL before logical design (wrong schema)
    \item Write queries before sample data exists (untestable)
    \item Approve a booking without validating against constraints (logical error)
\end{itemize}

Prerequisite enforcement (via SKILL.md) prevents these failures.

\subsection{Why Traceability Is Non-Negotiable}

In a graded academic project:
\begin{itemize}
    \item Grader must verify no invented features (no traceability = grade penalty)
    \item Student must justify every design choice (explicit matrix shows justification)
    \item Changes to requirement can be applied systematically (grep for entity, update all affected tables)
\end{itemize}

\section{Limitations \& Future Improvements}

\subsection{Current Limitations}
\begin{itemize}
    \item \textbf{Manual prerequisite checks}: Agent asks user to confirm override; not fully automated
    \item \textbf{No SQL validation}: DDL syntax checked by human, not machine
    \item \textbf{No conflict detection}: Booking overlaps checked in queries, not at constraint level
    \item \textbf{Mermaid rendering}: ERD exists but not rendered in PDF report
\end{itemize}

\subsection{Recommended Future Improvements}
\begin{enumerate}
    \item \textbf{Automated SQL parsing}: Pre-run DDL through SQL Server parser before submission
    \item \textbf{Multi-DBMS support}: Extend SKILL.md templates to PostgreSQL, MySQL, SQLite
    \item \textbf{Graphical ERD export}: Generate PNG/SVG from Mermaid for embedded in reports
    \item \textbf{Automated traceability check}: Machine-readable mapping (JSON) to verify coverage
    \item \textbf{Version control integration}: Auto-commit MEMORY.md after each step for audit trail
\end{enumerate}

\section{Conclusion}

The improved agent transforms the baseline 2-step demo into a comprehensive 7-step database design pipeline with state tracking, explicit traceability, validation, and prerequisite enforcement.

The markdown configuration files form a layered architecture:
\begin{enumerate}
    \item \textbf{AGENTS.md} provides static rules (workflow, naming, design principles) that all tools read automatically.
    \item \textbf{MEMORY.md} provides dynamic state (progress, decisions, questions) that enables cheap session resumption.
    \item \textbf{PROMPTS.md} provides human-facing prompt templates that enforce single-file, minimal-cost interactions.
    \item \textbf{SKILL.md} provides machine-facing step templates with prerequisite enforcement and done criteria.
\end{enumerate}

Every constraint traces directly to a teacher guideline from the reference project. The system supports four agent tools (Antigravity, OpenCode, Codex CLI, OpenClaw) without rule duplication, and enables auditable, cost-efficient database design work. All improvements are documented, reversible, and traceable to specific design weaknesses in the baseline agent.

\section*{Appendix: File Summary}

\begin{table}[h]
\centering
\begin{tabular}{|l|p{6cm}|}
\hline
\textbf{File} & \textbf{Purpose} \\
\hline
AGENTS.md & Global rules, workflow order, design principles \\
\hline
MEMORY.md & Project state, locked decisions, open questions, traceability \\
\hline
PROMPTS.md & Five templated prompts for different workflow scenarios \\
\hline
.opencode/skills/db-design-pipeline/SKILL.md & Detailed per-step templates, prerequisites, validation criteria \\
\hline
md/setup.md & Per-tool installation and launch instructions \\
\hline
md/pipeline.md & Project structure tree and pipeline overview \\
\hline
outputs/01--07-* & Seven deliverable artifacts (analysis, ERD, logical, validation, DDL, DML, queries) \\
\hline
\end{tabular}
\end{table}

\end{document}
